% paper/sections/10b_dccd_bootstrap.tex
% ═════════════════════════════════════════════════════════════════════════════
\subsection{DCCD パラメタの設計と適応制御}
\label{sec:dccd_params}
% ═════════════════════════════════════════════════════════════════════════════

§~\ref{sec:dccd_filter_theory} で導入した DCCD フィルタの理論的枠組みに基づき，
本節では散逸強度 $\varepsilon_{d,\max}$ の具体的な設計値と，
CLS 変数 $\psi$ に基づく適応制御アルゴリズムを与える．

\noindent\textbf{本節のスコープ．}
本稿では bulk 運動量移流の既定は UCCD6（§~\ref{sec:uccd6_def}）であり，
DCCD は (i) CLS $\psi$ 移流と再初期化圧縮項（§~\ref{sec:reinit}），
(ii) PPE RHS / Corrector のチェッカーボード抑制，
(iii) 曲率計算の高次微分安定化，の 3 用途に限定される．
以下の適応制御設計はこれらの用途における $\varepsilon_d$ と $S(\psi)$ の仕様である．

表~\ref{tab:dccd_usage_map} に，アルゴリズム全体で DCCD フィルタが担う 3 つの役割と
パラメータ設定を一覧する．

\begin{table}[htbp]
\centering
\caption{DCCD フィルタのパラメータ設計一覧}
\label{tab:dccd_usage_map}
\small
\begin{tabular}{llcl}
\hline
用途 & $\varepsilon_d$ & 適応制御 & 設計根拠 \\
\hline
CLS 移流（§~\ref{sec:advection_motivation}）
  & 0.05 & あり（$S(\psi)$） & 界面プロファイル保持 \\
PPE RHS / Corrector（§~\ref{sec:dccd_decoupling}）
  & $1/4$ & なし（固定） & チェッカーボード完全抑制 \\
曲率計算（§~\ref{sec:curvature}）
  & 0.05 & あり（$S(\psi)$） & 高次微分の精度保持 \\
\hline
\end{tabular}
\end{table}

% ─────────────────────────────────────────────────────────────────────────────
\subsubsection{散逸強度 \texorpdfstring{$\varepsilon_{d,\max}$}{ε\_d,max} の設計値}
\label{sec:dccd_eps_design}
% ─────────────────────────────────────────────────────────────────────────────

フィルタ伝達関数 $H(\xi;\varepsilon_d) = 1 - 4\varepsilon_d\sin^2(\xi/2)$ に対し，
低波数精度保存（波長 $\geq 6h$ で劣化 $\leq 5\%$）と
高周波抑制（ナイキスト成分 20\% 減衰）の同時要求から，
両制約が等号で一致する設計点が得られる
（スペクトル設計の詳細は §~\ref{sec:dccd_filter_theory} 参照）：
\begin{equation}
  \boxed{\varepsilon_{d,\max} = 0.05}
  \label{eq:eps_max}
\end{equation}
安定性制約 $\varepsilon_{d,\max} \leq 0.25$（式~\eqref{eq:dccd_filter_transfer}）を満足する．
より強い安定化が必要な場合は $\varepsilon_{d,\max} = 0.10$〜$0.20$ まで拡張可能だが，
低波数精度も同率で劣化する．

% ─────────────────────────────────────────────────────────────────────────────
\subsubsection{適応散逸制御則と実装アルゴリズム}
\label{sec:dccd_adaptive}
% ─────────────────────────────────────────────────────────────────────────────

\textbf{適応散逸強度．}
格子点 $i$ での散逸強度を CLS 変数 $\psi_i^n$（時刻 $t^n$）によって点別に制御する：
\begin{equation}
  \varepsilon_d^{(i)} \;=\; \varepsilon_{d,\max} \cdot S(\psi_i^n)
  \;=\; \varepsilon_{d,\max}\,(2\psi_i^n - 1)^2
  \label{eq:adaptive_eps}
\end{equation}
これにより：
\begin{itemize}
  \item $\psi_i \approx 0.5$（界面）：$\varepsilon_d^{(i)} \approx 0$，フィルタ無効，標準 CCD に退化
  \item $\psi_i \approx 0$ または $1$（バルク）：$\varepsilon_d^{(i)} \approx \varepsilon_{d,\max}$，
        最大散逸でノイズを積極的に排除
\end{itemize}

\begin{tcolorbox}[algbox, title={適応 Dissipative-CCD フィルタの実装アルゴリズム}]
\label{alg:dccd}
\textbf{入力：} スカラー場 $f$ の格子値 $\{f_i\}$，CLS 変数 $\{\psi_i^n\}$，散逸強度 $\varepsilon_{d,\max}$

\begin{enumerate}
  \item \textbf{CCD 求解：}
        標準 CCD ブロック Thomas アルゴリズム（§~\ref{sec:ccd_matrix}）を実行し，
        $\{f'_i\}$，$\{f''_i\}$ を得る．
  \item \textbf{点別散逸強度の計算：}
        各格子点 $i$ に対して
        \[
          \varepsilon_d^{(i)} = \varepsilon_{d,\max}\,(2\psi_i^n - 1)^2
        \]
        を計算する（計算コスト：$O(N)$，数値演算 1 回/点）．
  \item \textbf{1 次微分のフィルタ適用：}
        \[
          \tilde{f}'_i = f'_i + \varepsilon_d^{(i)}\,(f'_{i+1} - 2f'_i + f'_{i-1}),
          \quad i = 1, \ldots, N-2
        \]
        境界点 $i=0, N-1$：周期 BC では折り返しインデックスを使用；
        壁面・流出 BC では境界点をフィルタ適用前の CCD 値のまま保持する（§~\ref{sec:dccd_bc} 参照）．
  \item \textbf{2 次微分のフィルタ適用：}
        \[
          \tilde{f}''_i = f''_i + \varepsilon_d^{(i)}\,(f''_{i+1} - 2f''_i + f''_{i-1}),
          \quad i = 1, \ldots, N-2
        \]
  \item \textbf{後続計算への引き渡し：}
        フィルタ適用済みの $(\tilde{f}', \tilde{f}'')$ を CLS 移流，曲率計算，
        速度勾配評価に用いる．
        PPE 右辺・Corrector 発散のチェッカーボード抑制には，
        本適応制御ではなく $\varepsilon_d=1/4$ の固定フィルタ
        （§~\ref{sec:dccd_decoupling}）を用いる．
\end{enumerate}

\textbf{2 次元への拡張：}
2 次元計算では $x$ 方向と $y$ 方向に独立に 1 次元フィルタを適用する（次元分割法）：
\[
  \tilde{f}'_{x,ij} = f'_{x,ij} + \varepsilon_d^{(ij)}\,(f'_{x,i+1,j} - 2f'_{x,ij} + f'_{x,i-1,j}),
  \quad
  \tilde{f}'_{y,ij} = f'_{y,ij} + \varepsilon_d^{(ij)}\,(f'_{y,i,j+1} - 2f'_{y,ij} + f'_{y,i,j-1})
\]
$f''_{xx}, f''_{yy}, f''_{xy}$ に対しても同様に適用する．

\textbf{計算コスト：} ステップ 1（CCD，$O(N^2)$ for 2D）+ ステップ 2--4（フィルタ，$O(N^2)$）．
フィルタは CCD の後処理として追加コストが小さい．
また，$\tilde{f}'$ と $\tilde{f}''$ はフィルタ段階で\textbf{独立}に補正する——
すなわち $\tilde{f}''$ は $\tilde{f}'$ から再連立せずに計算する（計算コスト削減のための設計選択）．
\end{tcolorbox}

\noindent\textbf{BF 整合性への注意（P-5, \S\ref{sec:bf_p5_curvature}）}：
曲率平滑化を目的とする DCCD フィルタ適用は$\psi$・速度・$\kappa$ の
\textbf{帯制限}に限定し，圧力 $p$ および毛管圧 $p_\sigma = \sigma\kappa$ には
絶対に適用しない（\S\ref{sec:dccd_pressure_filter_prohibition}）．
片側だけのフィルタ適用は BF P-1 対称性を破壊する
（\S\ref{sec:dccd_pressure_filter_asymmetric}）．
本節の $\varepsilon_d^{(i)} = \varepsilon_{d,\max}(2\psi_i-1)^2$ 適応制御は
$S(\psi)$ が界面で消失するため界面帯制限の自然な実装となっており，
P-5 と整合する．

% ═════════════════════════════════════════════════════════════════════════════
\subsection{ブートストラップ処理：初期化シーケンス}
\label{sec:bootstrap}
% ═════════════════════════════════════════════════════════════════════════════

シミュレーション開始前に，格子--界面の循環参照を解消するため
以下の 6 段階ブートストラップを一度だけ実行する：
\begin{enumerate}[leftmargin=2em,noitemsep]
  \item \textbf{Step 0（一様格子仮 SDF）：} 等間隔格子 $\Delta x_0 = L/N$ 上で
        解析初期界面から仮の符号付距離 $\phi^{(0)}$ を評価する（幾何解析式または FMM）．
  \item \textbf{Step 1（非一様格子生成）：} Ridge-Eikonal 誘導 $\phi^{(0)}$ から
        界面適合非一様格子 $\{x_i\}, \{y_j\}$ を生成（§~\ref{sec:grid_2d}；
        密度比 $\alpha = \max h / \min h$）．
  \item \textbf{Step 2（CLS 初期化）：} 新しい非一様格子上で
        $\psi^0_{i,j} = \tfrac{1}{2}(1+\tanh(\phi^{(0)}_{i,j}/\varepsilon_c))$ を設定
        （$\varepsilon_c$ は CLS 厚さ，§~\ref{sec:cls_compression}）．
  \item \textbf{Step 3（CCD 演算子構築）：} 非一様格子重みから
        $D^{(1)}, D^{(2)}, M_{\mathrm{CCD}}, L_L, L_H$ の疎行列構造を一度だけ構築
        （格子更新 Mode 1 では以降不変）．
  \item \textbf{Step 3.5（非一様格子上 SDF 再初期化）：}
        Step 0 の $\phi^{(0)}$ は一様格子幾何解析由来であり，
        Step 1 で生成した非一様格子 $\{x_i\},\{y_j\}$ 上では
        $|\nabla\phi^{(0)}|=1$ が一般に保たれない．
        この差は Step 4 の Young--Laplace 圧力が依存する界面曲率
        $\kappa^0=\nabla\!\cdot\!(\nabla\phi^{(0)}/|\nabla\phi^{(0)}|)$
        の精度を律速し，かつ §~\ref{sec:field_extension} の HFE 精度要件
        （$\Ord{h^6}$）と整合しない．
        対策として，Step 3 の演算子で CLS 再初期化（§~\ref{sec:reinit}）を
        \textbf{1 サイクル}実行し
        $\phi^{(0)}\to\phi^{(0)}_\text{sdf}$ に高精度化する．
        \textbf{精度根拠}：広帯域 $\sigma_0 = 3 h_{\min}$ の 1 サイクルで
        $\bigl||\nabla\phi^{(0)}_\text{sdf}|-1\bigr|\le \Ord{h^4}$ が達成される
        （§~\ref{sec:reinit} の局所収束解析）．
        HFE 誤差伝播 $\varepsilon_\text{HFE} \lesssim \varepsilon_\text{sdf}\cdot|\nabla^k p|$
        ($k\le 6$) より局所 $\Ord{h^4}$ 誤差が界面曲率精度を律速するが，
        Step 4 Young--Laplace は時刻 0 の初期化のみであり，
        以後の時刻発展で誤差は減衰，
        \textbf{HFE $\Ord{h^6}$ 全体精度}は時間平均的に達成される
        （\S\ref{sec:verify_time} で数値検証）．
        以後の Step 4 では $\phi^{(0)}_\text{sdf}$ を用いる．
  \item \textbf{Step 4（Young--Laplace 初期圧力）：} 高精度化された
        $\phi^{(0)}_\text{sdf}$ から計算される初期曲率 $\kappa^0$ を用いて
        $p^0 = p_\infty + \sigma\kappa^0 H(\phi^{(0)}_\text{sdf})$（smoothed）で初期圧力場を設定し，
        Step 0 で測定した $\mathbf{u}^0=\mathbf{0}$ と整合させる．
\end{enumerate}
6 段階実行後，時刻ループ Step 1 から通常アルゴリズム（§\ref{sec:algo_flow}）が動作する．
詳細は付録~\ref{app:bootstrap}；格子の時間更新モード（Mode 1: 固定格子，Mode 2: 再生成）は
§~\ref{sec:grid_ale} を参照されたい．

% ═════════════════════════════════════════════════════════════════════════════
\subsection{タイムステップ制御}
\label{sec:algo_timestep}
% ═════════════════════════════════════════════════════════════════════════════

各タイムステップの開始時に，以下の2制約のうち最も厳しい条件で $\Delta t$ を決定する
（詳細な導出は §~\ref{sec:stability} 参照）：

\begin{equation}
  \Delta t = \min\!\left(
    \underbrace{
      C_{\mathrm{CFL}}
      \left[\max\!\left(\frac{|u|_{\max}}{\Delta x}
                       +\frac{|v|_{\max}}{\Delta y}\right)\right]^{-1}
    }_{\Delta t_{\mathrm{adv}}\text{：対流 CFL（式~\eqref{eq:dt_adv}）}},\;
    \underbrace{\sqrt{\frac{(\rho_l+\rho_g)\,h_{\min}^3}{2\pi\sigma}}}_{\Delta t_\sigma\text{：毛管波制約（式~\eqref{eq:dt_sigma}）}}
  \right)
  \label{eq:dt_full_algo}
\end{equation}

ここで $C_{\mathrm{CFL}} \le 0.5$ は安全率（デフォルト 0.45），$h_{\min} = \min(\Delta x, \Delta y)$ である．
表面張力なし（$\sigma=0$）の計算では毛管波制約を除外し，対流 CFL を用いる．
粘性項は Crank--Nicolson 法（無条件安定）で処理するため，等粘性系では粘性 CFL 制約は不要である．
ただし大粘性比（$\mu_l/\mu_g \gg 1$）の気液系では，クロス微分項 $\partial_y(\mu\partial_x v)$ の
陽的処理に起因する実質的な粘性制約が界面近傍で復活する場合がある（§~\ref{sec:stability} 参照）．

\noindent\textbf{§~\ref{sec:level_selection} の 3 項版との関係：}
式~\eqref{eq:dt_full_algo} の 2 項制約は等粘性・CN 有効領域での縮約表現であり，
一般の可変 Level 構成（§~\ref{sec:level_selection} 式~\eqref{eq:level_cfl_min}）では
$\Delta t_\text{visc}$（陽的粘性領域）も追加する 3 項 $\min$ が正しい．
物理波速度 $\Delta t_\text{wave}$ は非圧縮解法のため両式とも除外される．

\bigskip
以上の統合アルゴリズムの各構成要素が設計通りの精度を達成するか，
次章で系統的にコンポーネント単位の数学的検証を行う．
